\subsection{Альтернатива Фредгольма}

\subsubsection*{Четыре уравнения}

Основное уравнение и соответствующее однородное:
\[
  u(x) - \int\limits_{0}^{\ell} K(x,y)\, u(y)\, dy = f(x)
  \ \text{ в } \Lp{2}(0,\ell) ,
  \qquad\text{или}\qquad
  u - Au = f
  \tag{I}\label{eq:fa-main}
\]
\[
  u(x) - \int\limits_{0}^{\ell} K(x,y)\, u(y)\, dy = 0 ,
  \qquad\text{или}\qquad
  u - Au = 0
  \tag{$\mathrm{I}_0$}\label{eq:fa-hom}
\]

Чтобы выписать сопряжённое уравнение, найдём ядро оператора $A^*$:
\[
  (Au, v)
  = \int\limits_{0}^{\ell} \left[\ \int\limits_{0}^{\ell} K(x,y) u(y)\, dy
    \right] \overline{v(x)}\, dx
  = \int\limits_{0}^{\ell} u(y)\,
    \overline{\left[\ \int\limits_{0}^{\ell} \overline{K(x,y)}\, v(x)\, dx
    \right]} dy
  = (u, A^* v) ,
\]
то есть ядро сопряжённого оператора есть $\overline{K(y,x)}$:
\[
  (A^* v)(x) = \int\limits_{0}^{\ell} \overline{K(y,x)}\, v(y)\, dy .
\]
Отсюда ещё два уравнения:
\[
  v(x) - \int\limits_{0}^{\ell} \overline{K(y,x)}\, v(y)\, dy = g(x)
  \ \text{ в } \Lp{2}(D) ,
  \qquad\text{или}\qquad
  v - A^* v = g
  \tag{$\mathrm{I}^*$}\label{eq:fa-adj}
\]
\[
  v(x) - \int\limits_{0}^{\ell} \overline{K(y,x)}\, v(y)\, dy = 0 ,
  \qquad\text{или}\qquad
  v - A^* v = 0
  \tag{$\mathrm{I}_0^*$}\label{eq:fa-adj-hom}
\]

\begin{definition}
Говорят, что уравнение $Bu = c$ \textbf{фредгольмово разрешимо}, если из
единственности решения следует его существование, то есть из теоремы
единственности вытекает теорема существования.
\end{definition}

\begin{idea}
Формулировка непривычная, но смысл простой: обычно существование решения —
более трудный вопрос, чем единственность. Фредгольмова разрешимость говорит,
что для таких уравнений трудная половина достаётся даром: достаточно
проверить, что решений не больше одного.
\end{idea}

\subsubsection*{Теорема}

\begin{theorem}[альтернатива Фредгольма]\label{thm:fredholm}
Пусть $A$ — линейный компактный оператор в $H$. Тогда:
\begin{enumerate}[label=\arabic*), leftmargin=2em]
  \item уравнение~\eqref{eq:fa-main} разрешимо тогда и только тогда, когда
        $f$ ортогонален всем решениям сопряжённого однородного
        уравнения~\eqref{eq:fa-adj-hom}:
        \[
          f \in \big[\ker(I - A^*)\big]^{\perp} ;
        \]
  \item \emph{(альтернатива)} либо уравнение~\eqref{eq:fa-main} разрешимо
        единственным образом при любом $f \in H$ — и тогда оператор
        $(I-A)^{-1}$ определён на всём $H$ и ограничен, — либо однородное
        уравнение~\eqref{eq:fa-hom} имеет нетривиальные решения;
  \item уравнения~\eqref{eq:fa-hom} и~\eqref{eq:fa-adj-hom} имеют одинаковое
        число линейно независимых решений:
        \[
          \dim \ker(I - A) = \dim \ker(I - A^*) < \infty .
        \]
\end{enumerate}
\end{theorem}

\begin{remark}
Пункт~2 и есть фредгольмова разрешимость для линейных компактных операторов:
если уравнение~\eqref{eq:fa-main} может иметь не более одного решения, то оно
разрешимо при любой правой части.
\end{remark}

\subsubsection*{Доказательство: сведение к разобранным случаям}

Идея одна: компактный оператор расщепляется на конечномерную часть (к ней
применим \S4) и сжимающий остаток (к нему применим \S5).

\begin{proof}[Начало доказательства]
Оператор $A$ компактен, поэтому (\S2) его ядро раскладывается в ряд
\[
  K(x,y) = \sum_{i,j=1}^{\infty} k_{ij}\, \varphi_i(x)\,\overline{\psi_j(y)} ,
\]
сходящийся в $\Lp{2}(D\times D)$. Обозначим через $K_N$ частичную сумму
\[
  K_N(x,y) = \sum_{i,j=1}^{N} k_{ij}\, \varphi_i(x)\,\overline{\psi_j(y)}
\]
и положим: $A_1$ — оператор с ядром $K_N$, $A_2 = A - A_1$ — оператор с ядром
$K - K_N$. Тогда
\[
  A = A_1 + A_2 ,
  \qquad
  A_1 \ \text{— конечномерный},
  \qquad
  \norm{A_2} \le \norm{K - K_N} \xrightarrow[N\to\infty]{} 0 ,
\]
так что при достаточно большом $N$ выполнено $\norm{A_2} < 1$.

\inputfig{figures/fig-kernel-split}

\medskip
Перепишем уравнение~\eqref{eq:fa-main}:
\[
  u - A_1 u - A_2 u = f
  \qquad\Longleftrightarrow\qquad
  (I - A_2)\,u - A_1 u = f .
\]
По утверждению~\ref{prop:contraction} оператор $I - A_2$ обратим, поэтому
можно сделать замену
\[
  v := (I - A_2)\, u ,
  \qquad\text{то есть}\qquad
  u = (I - A_2)^{-1} v .
\]
Подставляя, получаем

\stepcounter{equation}
\begin{equation}\tag{1}\label{eq:fa-reduced}
  v - \underbrace{A_1 (I - A_2)^{-1}}_{=:\ B}\, v = f ,
  \qquad\text{то есть}\qquad
  v - Bv = f .
\end{equation}
Оператор $B$ \textbf{конечномерен}, поскольку его образ лежит в образе $A_1$,
а замена $u \leftrightarrow v$ взаимно однозначна. Значит
уравнения~\eqref{eq:fa-main} и~\eqref{eq:fa-reduced} эквивалентны, и к
последнему применима теорема~\ref{thm:deg-alternative}.
\end{proof}

\begin{warn}
На этом месте лекция обрывается. Остаётся довести перенос выводов с
уравнения~\eqref{eq:fa-reduced} обратно на~\eqref{eq:fa-main}: проверить, что
при замене $u \leftrightarrow v$ сохраняются число линейно независимых решений
однородных уравнений и условие ортогональности для правой части. Это,
по-видимому, начало следующей лекции.
\end{warn}

\begin{idea}[title={Почему расщепление устроено именно так}]
Два разобранных частных случая покрывают противоположные крайности:
\begin{itemize}[leftmargin=1.6em, topsep=2pt]
  \item \textbf{вырожденное ядро} (\S4) — оператор «большой», но действует в
        конечномерное пространство, и всё сводится к линейной алгебре;
  \item \textbf{сжимающий оператор} (\S5) — оператор может быть каким угодно
        по структуре, но мал по норме, и всё решается итерациями.
\end{itemize}
Компактность — это ровно то свойство, которое позволяет разложить оператор на
сумму таких двух: конечномерная часть забирает «большую» долю, а остаток
оказывается сколь угодно малым по норме. Если бы $A$ был просто ограниченным,
такого разложения не было бы, и альтернатива перестала бы работать.
\end{idea}
